\section{Conclusão}

\paragraph{}
A atividade permitiu aplicar diferentes técnicas de reconhecimento e
escalonamento de privilégios em sistemas Linux, utilizando tanto ferramentas
automatizadas quanto análise manual das informações obtidas.

\paragraph{}
Durante os testes foram identificadas e exploradas configurações inseguras
relacionadas a arquivos críticos do sistema, permissões inadequadas em
configurações do \texttt{sudo}, utilização indevida do grupo \texttt{docker} e
tarefas agendadas pelo \texttt{cron} executadas com privilégios elevados.

\paragraph{}
As ferramentas \texttt{LinEnum} e \texttt{PEASS} demonstraram sua utilidade na
fase de enumeração, auxiliando na identificação de potenciais vetores de ataque
e reduzindo significativamente o tempo necessário para o levantamento de
informações relevantes.

\paragraph{}
Os resultados obtidos evidenciam como erros de configuração aparentemente
simples podem permitir o comprometimento completo de um sistema Linux,
possibilitando a obtenção de privilégios administrativos por usuários não
autorizados.

\paragraph{}
O exercício reforçou a importância da correta configuração de permissões,
controle adequado de acesso a recursos privilegiados, revisão periódica de
tarefas automatizadas e aplicação do princípio do menor privilégio como medidas
essenciais para a segurança de ambientes Linux.
